\chapter{dbt Core: Transformasi, Dokumentasi, dan Lineage}
\label{modul:dbt-transformasi-dokumentasi-lineage}

\section{Identitas Modul}

\begin{identitasmodul}
\textbf{Sumber materi}    & Pertemuan 10 dan 11 \\
\textbf{CPMK}             & CPMK-092 dan CPMK-102 \\
\textbf{Minggu pelaksanaan} & Minggu ke-11 \\
\textbf{Durasi tatap muka} & 120 menit \\
\textbf{Estimasi tugas}   & 4--5 jam di luar sesi \\
\textbf{Moda kerja}       & Individual \\
\textbf{Perangkat}        & dbt Core dan PostgreSQL 17 \\
\textbf{Dataset}          & Staging NusaMart hasil pipeline ETL Modul 7 \\
\textbf{Skrip pendukung}  & Kerangka proyek dbt pada \berkas{dbt/modul-08/} \\
\textbf{Luaran}           & Proyek dbt yang berjalan, tangkapan graf lineage, dan DDL audit dimension \\
\end{identitasmodul}

\slideref{10 dan 11}

\section{Capaian dan Indikator Keberhasilan}

Mahasiswa mampu menyusun transformasi berlapis dari source ke staging dan mart,
mendokumentasikan model, membaca lineage, serta menjalankan pengujian dasar.
Target minimum modul adalah \perintah{dbt build} tanpa galat dan lineage utuh
dari source hingga mart.

\begin{capaian}
Pada akhir modul ini mahasiswa mampu:
\begin{enumerate}[leftmargin=1.4em]
  \item menjelaskan apa yang diperoleh ketika transformasi diperlakukan sebagai
        kode yang diversikan, bukan sebagai skrip yang dijalankan manual;
  \item mendefinisikan \perintah{source} dan membangun model berlapis dari
        staging ke mart dengan \perintah{ref};
  \item mendokumentasikan model dan kolom pada \berkas{schema.yml}, lalu
        membaca graf lineage yang dihasilkan \perintah{dbt docs generate};
  \item membangun \istilah{audit dimension} sehingga setiap baris fakta dapat
        ditelusuri ke batch pemuatan yang menghasilkannya;
  \item memasang empat generic data test dan menafsirkan kegagalannya.
\end{enumerate}
\end{capaian}

Modul ini menopang dua pertemuan sekaligus --- Pertemuan 10 tentang proses
ETL/ELT dan Pertemuan 11 tentang metadata, lineage, dan auditabilitas.
Keduanya bertemu pada satu gagasan: transformasi yang ditulis sebagai kode
menghasilkan metadata-nya sendiri sebagai efek samping.

\section{Prasyarat dan Persiapan}

\subsection{Prasyarat}

\begin{itemize}
  \item Pipeline dan schema \perintah{staging} hasil Modul 7 yang sudah terisi;
  \item SQL tingkat menengah: \perintah{CTE}, agregasi, dan join berlapis;
  \item Jinja tingkat dasar --- cukup memahami bahwa \perintah{\{\{ ... \}\}}
        adalah tempat yang diganti sebelum SQL dijalankan;
  \item konsep metadata teknis: nama objek, tipe kolom, dan asal-usul data;
  \item Git dasar, karena proyek dbt hanya bermakna bila diversikan.
\end{itemize}

\subsection{Pre-lab}

\begin{enumerate}[leftmargin=1.4em]
  \item Pasang dbt Core beserta adapter PostgreSQL:
\begin{lstlisting}[style=shellgaya]
python -m pip install dbt-core dbt-postgres
dbt --version
\end{lstlisting}
  \item Salin kerangka proyek yang disediakan asisten ke repositori Anda:
\begin{lstlisting}[style=shellgaya]
cp -r dbt/modul-08 modul-08-dbt/
cd modul-08-dbt
\end{lstlisting}
  \item Susun \berkas{\textasciitilde/.dbt/profiles.yml} dari contoh yang
        disediakan. Berkas ini berada di luar repositori --- lihat peringatan
        pada Bagian~A.
  \item Siapkan jawaban berikut pada \berkas{modul-08-dbt/pre-lab.md}:
  \begin{enumerate}[label=\alph*.,leftmargin=1.4em]
    \item Pada Modul 7 Anda menulis urutan pemuatan sendiri di dalam
          \berkas{pipeline.py}. Bagaimana Anda memastikan dimensi selalu dimuat
          sebelum fakta?
    \item Bila seorang rekan menambah satu model baru yang bergantung pada
          model Anda, bagaimana ia tahu urutan yang benar?
    \item Sebutkan satu pertanyaan tentang data yang \emph{tidak} dapat dijawab
          dengan membaca kode SQL saja.
  \end{enumerate}
\end{enumerate}

\section{Konsep Dasar}

\subsection{Transformasi sebagai Kode}

Sampai Modul 7, urutan pemuatan ditentukan oleh urutan pemanggilan fungsi di
dalam \berkas{pipeline.py}. Urutan itu benar, tetapi hanya diketahui oleh yang
menulisnya --- dan harus diperbarui manual setiap kali ada model baru.

dbt membalik hubungan itu. Anda menulis \emph{apa} yang bergantung pada apa,
lalu dbt menyusun urutannya sendiri.

\begin{konsep}{Empat hal yang diperoleh}
\begin{enumerate}[leftmargin=1.4em]
  \item \textbf{Urutan otomatis.} Ketergantungan dinyatakan lewat
        \perintah{ref}, dan dbt menurunkan urutan build dari situ.
  \item \textbf{Version control.} Transformasi menjadi berkas teks yang dapat
        di-\emph{diff}, ditinjau, dan dikembalikan.
  \item \textbf{Dokumentasi yang menyatu.} Deskripsi ditulis di sebelah
        modelnya, bukan pada dokumen terpisah yang cepat basi.
  \item \textbf{Lineage gratis.} Graf ketergantungan adalah hasil sampingan
        dari \perintah{ref} --- tidak ada pekerjaan tambahan untuk
        memperolehnya.
\end{enumerate}
\end{konsep}

Butir keempat inilah alasan RANCANGAN praktikum ini menolak perkakas lineage
tersendiri. Lineage yang harus dipelihara terpisah akan selalu tertinggal dari
kodenya; lineage yang \emph{diturunkan} dari kode tidak bisa tertinggal.

\subsection{Sources, Models, dan Materialization}

Dua fungsi Jinja membangun seluruh graf, dan perbedaannya menentukan bentuk
lineage Anda.

\begin{center}
\small
\begin{tabular}{@{}p{0.22\textwidth}>{\raggedright\arraybackslash}p{0.34\textwidth}>{\raggedright\arraybackslash}p{0.36\textwidth}@{}}
\toprule
\textbf{Fungsi} & \textbf{Menunjuk} & \textbf{Dipakai untuk} \\
\midrule
\perintah{source()} & Tabel yang \emph{tidak} dibangun dbt &
  Tabel \perintah{staging} hasil pipeline Modul 7 \\
\perintah{ref()} & Model lain di dalam proyek dbt &
  Seluruh ketergantungan antarmodel \\
\bottomrule
\end{tabular}
\end{center}

\begin{peringatan}
Menulis nama tabel secara langsung --- \perintah{FROM staging.penjualan} alih-alih
\perintah{FROM \{\{ source('nusamart', 'penjualan') \}\}} --- tetap
menghasilkan SQL yang berjalan. Yang hilang adalah ketergantungannya: model itu
tidak muncul pada graf lineage, dan dbt tidak tahu ia harus dibangun sesudah
apa. Kesalahan ini tidak menimbulkan galat, hanya graf yang bohong.
\end{peringatan}

\istilah{Materialization} menentukan wujud model di basis data.

\begin{center}
\small
\begin{tabular}{@{}p{0.20\textwidth}>{\raggedright\arraybackslash}p{0.36\textwidth}>{\raggedright\arraybackslash}p{0.36\textwidth}@{}}
\toprule
\textbf{Materialization} & \textbf{Wujudnya} & \textbf{Dipakai untuk} \\
\midrule
\perintah{view} & \perintah{CREATE VIEW} & Model staging: murah, selalu segar \\
\perintah{table} & \perintah{CREATE TABLE} & Model mart: mahal dibangun,
  murah dibaca \\
\perintah{incremental} & Menambah baris baru saja & Fakta besar yang tumbuh \\
\perintah{ephemeral} & Disisipkan sebagai CTE & Perantara yang tidak perlu ada
  di basis data \\
\bottomrule
\end{tabular}
\end{center}

Aturan praktis yang dipakai modul ini: \emph{staging} sebagai
\perintah{view}, \emph{mart} sebagai \perintah{table}. Pertukarannya sama
dengan yang Anda ukur pada Modul 6 --- ruang dan waktu pembangunan ditukar
dengan waktu pembacaan.

\subsection{Dokumentasi dan Graf Lineage}

Deskripsi model dan kolom ditulis pada berkas \berkas{schema.yml} yang berada
di direktori yang sama dengan modelnya. Kedekatan ini disengaja: dokumentasi
yang berjauhan dari kodenya akan basi dalam hitungan minggu.

\perintah{dbt docs generate} menggabungkan tiga hal --- deskripsi yang Anda
tulis, katalog basis data yang dibaca dbt, dan graf ketergantungan dari
\perintah{ref} --- menjadi satu situs statis. Graf lineage yang muncul di
situ bukan gambar yang Anda buat, melainkan hasil pembacaan kode Anda sendiri.

\begin{konsep}{Dua arah membaca lineage}
\begin{itemize}
  \item \textbf{Ke hulu} (\emph{upstream}): angka ini berasal dari mana?
        Dipakai saat menelusuri angka yang mencurigakan.
  \item \textbf{Ke hilir} (\emph{downstream}): kalau saya mengubah model ini,
        apa saja yang ikut rusak? Dipakai sebelum mengubah apa pun.
\end{itemize}
\end{konsep}

\subsection{Audit Dimension}

Tabel \perintah{gudang.audit\_muat} yang Anda isi pada Modul 7 mencatat
\emph{jalannya} pipeline. \istilah{Audit dimension} melangkah lebih jauh: ia
menghubungkan catatan itu ke \emph{setiap baris fakta}, sehingga pertanyaan
``baris ini dimuat kapan, dari sumber mana, oleh pipeline versi berapa?''
dapat dijawab dengan satu join.

\begin{center}
\small
\begin{tabular}{@{}p{0.26\textwidth}>{\raggedright\arraybackslash}p{0.64\textwidth}@{}}
\toprule
\textbf{Kolom} & \textbf{Gunanya} \\
\midrule
\perintah{batch\_id} & Menghubungkan ke catatan \perintah{audit\_muat} \\
\perintah{waktu\_muat} & Kapan baris ini masuk gudang \\
\perintah{sumber} & Sistem operasional asalnya \\
\perintah{versi\_pipeline} & Kode versi berapa yang menghasilkannya \\
\perintah{status\_mutu} & Lulus, lulus dengan catatan, atau dikarantina \\
\bottomrule
\end{tabular}
\end{center}

Kegunaannya terasa ketika sesuatu salah. Bila ditemukan bahwa pipeline versi
tertentu mengandung cacat, audit dimension menjawab dengan tepat baris fakta
mana saja yang perlu dimuat ulang --- tanpanya, satu-satunya pilihan aman
adalah memuat ulang semuanya.

\subsection{Pengujian Data Dasar}

dbt menyediakan empat uji baku yang cukup untuk menangkap sebagian besar
kerusakan. Keempatnya ditulis sebagai beberapa baris YAML, bukan sebagai kode.

\begin{center}
\small
\begin{tabular}{@{}p{0.22\textwidth}>{\raggedright\arraybackslash}p{0.32\textwidth}>{\raggedright\arraybackslash}p{0.38\textwidth}@{}}
\toprule
\textbf{Uji} & \textbf{Memeriksa} & \textbf{Contoh pada NusaMart} \\
\midrule
\perintah{not\_null} & Kolom tidak kosong &
  Kunci dimensi pada fakta \\
\perintah{unique} & Nilai tidak berulang &
  \perintah{produk\_key} pada dimensi \\
\perintah{relationships} & Setiap nilai ada di tabel rujukan &
  Fakta menunjuk dimensi yang benar-benar ada \\
\perintah{accepted\_values} & Nilai berada dalam daftar &
  \perintah{status} hanya berisi nilai yang dikenal \\
\bottomrule
\end{tabular}
\end{center}

\begin{catatan}
Uji \perintah{relationships} adalah versi dbt dari \perintah{FOREIGN KEY}, dan
keduanya tidak saling menggantikan. Constraint basis data menolak baris salah
\emph{saat dimasukkan}; uji dbt melaporkan baris salah \emph{setelah}
transformasi berjalan --- termasuk pada model yang berupa view, yang tidak
dapat diberi constraint sama sekali.
\end{catatan}

Modul ini memakai keempat uji sebagai pengenalan. Perancangan aturan kualitas
yang sesungguhnya --- termasuk custom test, ambang, dan karantina --- menjadi
bahan Modul 10.

\section{Studi Kasus dan Protokol Praktikum}

Mahasiswa melengkapi kerangka dbt yang telah disediakan. Fokus sesi adalah
memahami struktur proyek dan hubungan transformasi, bukan memulai seluruh
proyek dari nol.

Proyek dbt membangun ulang jalur \perintah{staging} sampai \perintah{mart}
sebagai kode, dengan keluaran pada schema \perintah{dbt\_mart} yang
\emph{terpisah} dari \perintah{gudang} bikinan tangan. Pemisahan ini
disengaja: pada akhir sesi, angka keduanya dibandingkan, dan keduanya harus
sama.

\begin{center}
\small
\begin{tabular}{@{}p{0.28\textwidth}p{0.16\textwidth}>{\raggedright\arraybackslash}p{0.48\textwidth}@{}}
\toprule
\textbf{Model} & \textbf{Lapis} & \textbf{Keadaan pada kerangka} \\
\midrule
\perintah{stg\_penjualan} & staging & Disediakan lengkap sebagai contoh \\
\perintah{stg\_produk}    & staging & Disediakan lengkap \\
\perintah{stg\_toko}      & staging & \textbf{Dilengkapi mahasiswa} \\
\perintah{dim\_audit}     & mart    & Disediakan lengkap \\
\perintah{dim\_produk}    & mart    & Disediakan lengkap sebagai contoh SCD-2 \\
\perintah{fct\_penjualan} & mart    & \textbf{Dilengkapi mahasiswa} \\
\bottomrule
\end{tabular}
\end{center}

\begin{catatanasisten}
Kerangka sengaja setengah jadi. Dua model yang dikosongkan dipilih supaya
mahasiswa harus memakai \perintah{source} sekali dan \perintah{ref} beberapa
kali --- keduanya wajib dipahami agar graf lineage terbentuk utuh.

Bila sesi berjalan lambat, prioritaskan Bagian B sampai D. Bagian E dapat
dipindahkan ke tugas luar laboratorium tanpa merusak alur modul, karena
pengujian dibahas lebih dalam pada Modul 10.
\end{catatanasisten}

\section{Alur Praktikum 120 Menit}

\begin{alurwaktu}
0--15 & Demonstrasi source, ref, build, dan docs. \\
15--95 & Kerja terbimbing: staging, mart, dokumentasi, lineage, audit, dan test. \\
95--110 & Checkpoint dbt build serta graf lineage. \\
110--120 & Penjelasan tugas penelusuran satu angka. \\
\end{alurwaktu}

\section{Prosedur Praktikum Terbimbing}

\subsection{Bagian A: Menghubungkan dbt ke PostgreSQL}

\langkah{A.1 Menyusun profil koneksi}{10 menit}

\begin{lstlisting}[style=yamlgaya]
# ~/.dbt/profiles.yml  -- DI LUAR repositori
nusamart:
  target: dev
  outputs:
    dev:
      type: postgres
      host: localhost
      port: 5434
      user: praktikum
      password: "{{ env_var('DW_PASSWORD') }}"
      dbname: nusamart_dw
      schema: dbt_mart
      threads: 4
\end{lstlisting}

\begin{peringatan}
Berkas \berkas{profiles.yml} memuat kredensial dan karena itu berada di
\berkas{\textasciitilde/.dbt/}, \textbf{bukan} di dalam repositori. Bentuk
\perintah{env\_var} membaca kata sandi dari lingkungan, sehingga berkasnya
sendiri tetap aman dibagikan. Repositori Anda hanya memuat
\berkas{profiles.yml.example} tanpa nilai sebenarnya.
\end{peringatan}

\langkah{A.2 Menguji koneksi}{5 menit}

\begin{lstlisting}[style=shellgaya]
export DW_PASSWORD='...'
dbt debug          # memeriksa profil, koneksi, dan berkas proyek
dbt deps           # memasang paket bila ada
\end{lstlisting}

Seluruh baris keluaran \perintah{dbt debug} harus bertanda \perintah{OK}.
Jangan lanjut sebelum itu tercapai --- galat koneksi yang muncul di tengah
pembangunan model jauh lebih sulit dibaca.

\subsection{Bagian B: Mendefinisikan Source dan Staging}

\langkah{B.1 Membaca definisi source}{10 menit}

\begin{lstlisting}[style=yamlgaya]
# models/sources.yml
version: 2

sources:
  - name: nusamart
    description: >
      Tabel staging hasil pipeline ETL Modul 7. Berisi salinan tiga sumber
      operasional yang sudah distandardisasi satuan dan kodenya.
    database: nusamart_dw
    schema: staging
    tables:
      - name: penjualan_gabungan
        description: Satu baris per produk per baris struk, dari tiga sumber.
        columns:
          - name: transaksi_id
            description: Nomor transaksi pada sistem sumber.
      - name: produk_gabungan
        description: Produk dari tiga sumber, berat sudah dalam gram.
\end{lstlisting}

Deskripsi bukan hiasan. Ia muncul pada situs dokumentasi, dan pada Bagian D
Anda akan membacanya kembali sebagai pengguna, bukan sebagai penulis.

\langkah{B.2 Melengkapi model staging}{15 menit}

Model \perintah{stg\_penjualan} disediakan lengkap sebagai contoh. Perhatikan
tiga hal: \perintah{config} materialization, pemakaian \perintah{source}, dan
penamaan kolom yang diseragamkan.

\begin{lstlisting}[style=sqlgaya]
-- models/staging/stg_penjualan.sql
{{ config(materialized='view') }}

WITH sumber AS (
    SELECT * FROM {{ source('nusamart', 'penjualan_gabungan') }}
)
SELECT transaksi_id,
       nomor_struk,
       produk_id,
       toko_id,
       pelanggan_master_id AS pelanggan_id,
       waktu_transaksi::DATE AS tanggal_transaksi,
       kuantitas,
       harga_satuan,
       diskon,
       subtotal,
       sumber        AS sistem_asal,
       batch_id
FROM   sumber
WHERE  status <> 'BATAL'
\end{lstlisting}

Lengkapi \perintah{stg\_toko} dengan pola yang sama, lalu bangun:

\begin{lstlisting}[style=shellgaya]
dbt run --select staging
\end{lstlisting}

\subsection{Bagian C: Membangun Model Mart}

\langkah{C.1 Membaca model mart contoh}{10 menit}

\begin{lstlisting}[style=sqlgaya]
-- models/mart/dim_produk.sql
{{ config(materialized='table') }}

SELECT produk_id,
       sku,
       nama_produk,
       kategori,
       berat_gram,
       mulai_berlaku,
       selesai_berlaku,
       baris_kini
FROM   {{ ref('stg_produk') }}
\end{lstlisting}

Perhatikan \perintah{ref}, bukan nama tabel. Inilah satu-satunya cara dbt
mengetahui bahwa model ini harus dibangun \emph{sesudah}
\perintah{stg\_produk}.

\langkah{C.2 Melengkapi fct\_penjualan}{20 menit}

Model ini sengaja dikosongkan. Ia harus memakai \perintah{ref} ke seluruh
model dimensi, dan menerapkan dua aturan yang sudah Anda pegang sejak Modul 3
dan Modul 5.

\begin{lstlisting}[style=sqlgaya]
-- models/mart/fct_penjualan.sql
{{ config(materialized='table') }}

SELECT f.transaksi_id,
       f.nomor_struk,
       COALESCE(p.produk_key, -1)   AS produk_key,   -- baris unknown
       COALESCE(t.toko_key, -1)     AS toko_key,
       a.audit_key,
       f.tanggal_transaksi,
       f.kuantitas,
       f.harga_satuan,
       f.diskon,
       f.subtotal
FROM       {{ ref('stg_penjualan') }} f
-- LEFT JOIN, dan pencarian versi memakai rentang tanggal transaksi
LEFT JOIN  {{ ref('dim_produk') }} p
       ON  p.produk_id = f.produk_id
       AND f.tanggal_transaksi >= p.mulai_berlaku
       AND f.tanggal_transaksi <  p.selesai_berlaku
LEFT JOIN  {{ ref('dim_toko') }} t ON t.toko_id = f.toko_id
LEFT JOIN  {{ ref('dim_audit') }} a ON a.batch_id = f.batch_id
\end{lstlisting}

\begin{lstlisting}[style=shellgaya]
dbt build          # run + test, mengikuti urutan graf
\end{lstlisting}

\langkah{C.3 Merekonsiliasi dengan gudang bikinan tangan}{5 menit}

\begin{lstlisting}[style=sqlgaya]
SELECT (SELECT SUM(subtotal) FROM dbt_mart.fct_penjualan)   AS versi_dbt,
       (SELECT SUM(subtotal) FROM gudang.fakta_penjualan)   AS versi_tangan,
       (SELECT SUM(subtotal) FROM dbt_mart.fct_penjualan)
     - (SELECT SUM(subtotal) FROM gudang.fakta_penjualan)   AS selisih;
\end{lstlisting}

Keduanya membangun hal yang sama dengan cara berbeda, sehingga angkanya harus
sama. Selisih yang muncul hampir selalu berasal dari penyaringan status atau
penanganan baris unknown yang tidak konsisten antara kedua jalur --- dan itu
temuan yang layak dicatat.

\subsection{Bagian D: Mendokumentasikan dan Membaca Lineage}

\langkah{D.1 Melengkapi schema.yml}{10 menit}

\begin{lstlisting}[style=yamlgaya]
# models/mart/schema.yml
version: 2

models:
  - name: fct_penjualan
    description: >
      Fakta penjualan, satu baris per produk per baris struk. Grain sama
      dengan gudang.fakta_penjualan yang dibangun manual pada Modul 2.
    columns:
      - name: produk_key
        description: >
          Kunci dimensi produk yang BERLAKU saat transaksi terjadi, bukan
          versi terkini. Bernilai -1 bila produk tidak dikenali.
      - name: audit_key
        description: Kunci batch pemuatan yang menghasilkan baris ini.
\end{lstlisting}

\langkah{D.2 Membangun dan membaca graf lineage}{10 menit}

\begin{lstlisting}[style=shellgaya]
dbt docs generate
dbt docs serve --port 8080
\end{lstlisting}

Buka graf lineage, lalu simpan tangkapan layarnya sebagai
\berkas{lineage.png}. Graf harus memperlihatkan jalur utuh dari
\perintah{source} sampai \perintah{fct\_penjualan}.

\begin{peringatan}
Model yang muncul \emph{terputus} dari graf --- tidak memiliki panah masuk ---
hampir selalu menandakan nama tabel ditulis langsung, bukan lewat
\perintah{source} atau \perintah{ref}. Perbaiki sebelum melanjutkan; graf yang
bohong lebih berbahaya daripada tidak ada graf sama sekali.
\end{peringatan}

Telusuri satu model ke hulu dan satu ke hilir, lalu catat pada
\berkas{catatan-lineage.md}: model apa yang akan ikut rusak bila
\perintah{stg\_penjualan} diubah?

\subsection{Bagian E: Menambah Audit dan Pengujian Dasar}

\langkah{E.1 Membaca audit dimension}{5 menit}

\begin{lstlisting}[style=sqlgaya]
-- models/mart/dim_audit.sql
{{ config(materialized='table') }}

SELECT ROW_NUMBER() OVER (ORDER BY batch_id, sumber) AS audit_key,
       batch_id,
       sumber,
       MIN(mulai)   AS waktu_muat,
       SUM(baris_dibaca)  AS baris_dibaca,
       SUM(baris_dimuat)  AS baris_dimuat,
       CASE WHEN SUM(baris_dibaca) = SUM(baris_dimuat) THEN 'lengkap'
            ELSE 'ada selisih' END AS status_mutu
FROM   {{ source('nusamart', 'audit_muat') }}
GROUP  BY batch_id, sumber
\end{lstlisting}

\langkah{E.2 Memasang empat uji baku}{10 menit}

\begin{lstlisting}[style=yamlgaya]
models:
  - name: fct_penjualan
    columns:
      - name: transaksi_id
        tests:
          - not_null
      - name: produk_key
        tests:
          - not_null
          - relationships:
              to: ref('dim_produk')
              field: produk_key
  - name: dim_produk
    columns:
      - name: produk_key
        tests:
          - unique
          - not_null
\end{lstlisting}

\begin{lstlisting}[style=shellgaya]
dbt test
dbt test --store-failures     # baris yang gagal disimpan untuk diperiksa
\end{lstlisting}

\begin{catatan}
Uji yang gagal bukan kegagalan modul. Yang dinilai adalah kemampuan
\emph{menafsirkan} kegagalannya: baris mana yang gagal, mengapa, dan apakah
yang harus diperbaiki adalah datanya, transformasinya, atau justru ujinya ---
karena kadang aturan yang ditulis memang terlalu ketat.
\end{catatan}

\begin{luaran}
\berkas{modul-08-dbt/} berisi proyek dbt lengkap tanpa kredensial, keluaran
\perintah{dbt build}, \berkas{lineage.png}, \berkas{catatan-lineage.md}, dan
DDL audit dimension.
\end{luaran}

\begin{checkpoint}
Asisten memeriksa butir berikut, sebagian dari keluaran \perintah{dbt} dan
sebagian dari \berkas{sql/modul-08/check.sql}:
\begin{ceklis}
  \item \perintah{dbt debug} seluruhnya \perintah{OK};
  \item \perintah{dbt build} berjalan tanpa galat;
  \item seluruh model memakai \perintah{source} atau \perintah{ref}, tidak ada
        nama tabel yang ditulis langsung;
  \item graf lineage utuh dari source sampai \perintah{fct\_penjualan}, tanpa
        model yang terputus;
  \item angka \perintah{dbt\_mart.fct\_penjualan} sama dengan
        \perintah{gudang.fakta\_penjualan};
  \item audit dimension terisi otomatis dan setiap baris fakta memiliki
        \perintah{audit\_key};
  \item empat generic test terpasang dan hasilnya dilaporkan;
  \item \berkas{profiles.yml} tidak ada di dalam repositori.
\end{ceklis}
\end{checkpoint}

\section{Latihan Mandiri di Kelas}

\latihan{Memprediksi dampak perubahan pada DAG}

Anda hendak menambahkan kolom \perintah{margin} pada \perintah{stg\_penjualan}.

\begin{enumerate}[leftmargin=1.4em]
  \item Sebelum melihat graf, tuliskan model apa saja yang menurut Anda ikut
        terpengaruh.
  \item Buktikan dengan dbt:
\begin{lstlisting}[style=shellgaya]
dbt ls --select stg_penjualan+      # ke hilir
dbt ls --select +fct_penjualan      # ke hulu
\end{lstlisting}
  \item Bandingkan prediksi dengan hasilnya, lalu jelaskan selisihnya.
  \item Berapa lama perkiraan Anda untuk menjawab pertanyaan yang sama pada
        pipeline Python Modul 7, yang tidak memiliki graf?
\end{enumerate}

\latihan{Membuat satu model rusak dengan sengaja}

Pada model \perintah{fct\_penjualan}, ganti pemanggilan
\perintah{ref('dim\_produk')} dengan nama tabel yang ditulis langsung, yaitu
\perintah{dbt\_mart.dim\_produk}.

\begin{enumerate}[leftmargin=1.4em]
  \item Jalankan \perintah{dbt build}. Apakah ia berhasil?
  \item Jalankan \perintah{dbt docs generate}. Apa yang berubah pada graf?
  \item Jelaskan mengapa kesalahan semacam ini lebih berbahaya daripada
        kesalahan yang membuat build gagal.
  \item Kembalikan ke bentuk semula.
\end{enumerate}

\section{Tugas Individu}

\subsection{Pekerjaan Wajib}
Telusuri satu angka dari mart atau dashboard sampai ke baris sumber dan tulis
jalur lineage tersebut secara naratif.

Kumpulkan \berkas{telusur-angka.md} yang memuat:

\begin{enumerate}[leftmargin=1.4em]
  \item satu angka tertentu yang Anda pilih --- misalnya nilai penjualan
        kategori tertentu pada bulan tertentu --- beserta nilainya;
  \item jalur lineage-nya secara naratif, dari model mart mundur sampai baris
        pada sistem sumber, menyebut setiap model yang dilewati;
  \item query yang Anda jalankan pada setiap langkah untuk membuktikan bahwa
        angka itu memang berasal dari sana;
  \item transformasi apa saja yang dialami angka tersebut sepanjang jalan ---
        penyaringan, konversi satuan, agregasi;
  \item batch pemuatan mana yang menghasilkannya, dibaca dari audit dimension.
\end{enumerate}

\begin{catatan}
Butir kelima adalah alasan audit dimension dibangun. Tanpanya, penelusuran
berhenti pada ``baris ini ada di staging'' --- dan pertanyaan ``masuk kapan,
dari jalan pipeline yang mana'' tidak dapat dijawab sama sekali.
\end{catatan}

\subsection{Pertanyaan Analisis}

\begin{enumerate}[leftmargin=1.4em]
  \item Seorang rekan mengganti \perintah{ref} dengan nama tabel langsung agar
        modelnya dapat dijalankan sendiri tanpa membangun ulang seluruh
        proyek. Sebutkan apa yang diperolehnya dan apa yang hilang, lalu
        berikan pendirian Anda.
  \item Uji \perintah{relationships} Anda lulus, tetapi baris fakta yang
        menunjuk baris unknown berjumlah besar. Apakah data Anda sehat?
        Jelaskan mengapa kedua hal itu tidak bertentangan, dan uji tambahan apa
        yang perlu ditulis.
  \item Dokumentasi dbt dibangun ulang setiap kali \perintah{docs generate}
        dijalankan. Sebutkan satu jenis pengetahuan tentang data yang
        \emph{tidak} akan pernah muncul di sana meskipun seluruh deskripsi
        terisi, dan di mana pengetahuan itu seharusnya disimpan.
\end{enumerate}

\section{Format Pengumpulan dan Checklist}

Kumpulkan proyek dbt tanpa kredensial, keluaran \perintah{dbt build}, tangkapan
layar lineage, DDL audit dimension, dan narasi penelusuran angka.

Seluruh berkas diletakkan pada \berkas{modul-08-dbt/}.

\begin{ceklis}
  \item \berkas{pre-lab.md}
  \item Proyek dbt lengkap: \berkas{dbt\_project.yml}, \berkas{models/},
        \berkas{macros/}
  \item \berkas{profiles.yml.example} --- \textbf{tanpa} kredensial
  \item \berkas{keluaran-dbt-build.txt}
  \item \berkas{lineage.png} dan \berkas{catatan-lineage.md}
  \item \berkas{telusur-angka.md}
  \item \berkas{jawaban-analisis.md}
  \item Riwayat commit memperlihatkan model ditambahkan bertahap, bukan
        sekaligus di akhir.
\end{ceklis}

\section{Troubleshooting}

\begin{table}[htbp]
\centering
\small
\caption{Masalah yang lazim muncul pada Modul 8 dan penanganannya}
\label{tab:m8-troubleshooting}
\begin{tabular}{@{}p{0.30\textwidth}>{\raggedright\arraybackslash}p{0.62\textwidth}@{}}
\toprule
\textbf{Gejala} & \textbf{Penanganan} \\
\midrule
\perintah{Could not find profile named 'nusamart'} &
  Nama profil pada \berkas{dbt\_project.yml} berbeda dengan kunci teratas
  \berkas{profiles.yml}. Keduanya harus sama persis. \\
\perintah{Env var required but not provided: 'DW\_PASSWORD'} &
  Variabel lingkungan belum diekspor pada sesi terminal ini. Jalankan
  \perintah{export} lebih dahulu. \\
\perintah{Compilation Error ... depends on a node named ... which was not
  found} &
  Nama model pada \perintah{ref} salah ketik, atau berkasnya berada di luar
  direktori \berkas{models/}. Nama yang dipakai adalah nama berkas tanpa
  \berkas{.sql}. \\
\perintah{Database Error ... relation "staging.xxx" does not exist} &
  Schema atau nama tabel pada \berkas{sources.yml} keliru, atau pipeline
  Modul 7 belum dijalankan. \\
Model muncul terputus pada graf &
  Nama tabel ditulis langsung, bukan lewat \perintah{source} atau
  \perintah{ref}. Build tetap berhasil --- inilah yang membuatnya berbahaya. \\
\perintah{dbt docs serve} menampilkan halaman kosong &
  \perintah{dbt docs generate} belum dijalankan, atau dijalankan dari
  direktori yang bukan akar proyek. \\
Uji \perintah{unique} gagal pada dimensi SCD-2 &
  Wajar dan benar: satu natural key memang punya beberapa versi. Uji
  \perintah{unique} dipasang pada surrogate key, bukan natural key. \\
Uji \perintah{relationships} gagal untuk kunci $-1$ &
  Baris unknown belum ada pada model dimensi versi dbt. Tambahkan, sebagaimana
  Modul 5. \\
Angka dbt berbeda dengan gudang bikinan tangan &
  Paling sering karena penyaringan status atau penanganan unknown tidak
  konsisten antara kedua jalur. Bandingkan klausa \perintah{WHERE}
  keduanya. \\
\bottomrule
\end{tabular}
\end{table}

\section{Ringkasan}

Modul ini memindahkan transformasi dari skrip yang dijalankan manual menjadi
kode yang menyusun urutannya sendiri, mendokumentasikan dirinya, dan
menghasilkan graf ketergantungannya sebagai efek samping.

Tiga gagasan perlu dibawa ke Modul 9 dan 10. \emph{Pertama}, ketergantungan
yang dinyatakan lewat \perintah{ref} adalah satu-satunya sumber kebenaran graf
lineage --- menuliskan nama tabel secara langsung tetap menghasilkan SQL yang
berjalan, tetapi graf yang bohong. \emph{Kedua}, metadata teknis paling
berguna ketika melekat pada barisnya: audit dimension mengubah pertanyaan
``baris ini dari mana'' dari penyelidikan menjadi satu join. \emph{Ketiga},
uji baku dbt menangkap kerusakan struktural, tetapi tidak menangkap angka yang
salah secara halus --- untuk itu diperlukan aturan yang dirancang khusus, dan
itulah Modul 10.

Modul 9 memakai schema \perintah{mart} yang kini memiliki dua penghuni ---
hasil kerja tangan dan hasil kerja dbt --- untuk membangun dashboard yang
menjawab pertanyaan bisnis yang Anda rumuskan pada Modul 1. Lingkaran
praktikum ditutup di sana.
